\section{Background modes and strict envelopes}
\label{sec:envelopes}

Let
\[
 M=S^2\times\mathbb T^2,\qquad
 g_0=dp^2+\sin^2p\,d\theta^2+dz^2+d{z'}^2,
\]
where the sphere has radius one and both circles have length $2\pi$.
Fix a coordinate cube $P$ whose closure lies in $0<p<\pi/2$, and a
smaller coordinate cube $D\Subset P$. Throughout the proof, integration
against $dx$ in this chart means coordinate Lebesgue measure.

\begin{lemma}[Background modes]\label{env:background-mode}
For every integer $k\ge1$, set
\[
 \Lambda_k=2k^2+k,\qquad
 u_k^0=\sin^kp\cos(k\phi),\qquad \phi=\theta+z.
\]
Then $u_k^0$ is a smooth real eigenfunction on $M$ with
$(\Delta_{g_0}+\Lambda_k)u_k^0=0$. Set $f_0=\log\sin p$ on a
neighborhood of $\overline P$. There is a constant $c>0$, independent
of $k$, such that
\begin{equation}\label{env:background-gradient}
 |du_k^0|_{g_0}\ge ck e^{kf_0}\quad\hbox{on }\overline P.
\end{equation}
\end{lemma}

\begin{proof}
In Cartesian coordinates on the sphere, $u_k^0$ is the real part of
$(X_1+iX_2)^ke^{ikz}$. The polynomial $(X_1+iX_2)^k$ is homogeneous
and harmonic, so its restriction to the unit sphere has eigenvalue
$k(k+1)$; the $z$ factor contributes $k^2$. This proves smoothness and
the eigenfunction equation. The covectors $df_0$ and $d\phi$ are
orthogonal, with squared norms $\cot^2p$ and $\csc^2p+1$. Hence
\begin{align}
 |du_k^0|_{g_0}^2
 &=k^2e^{2kf_0}\bigl(\cot^2p\cos^2(k\phi)
                  +(\csc^2p+1)\sin^2(k\phi)\bigr)\notag\\
 &=k^2\bigl(\sin^{2k}p+\sin^{2k-2}p-2(u_k^0)^2\bigr).
 \label{env:background-gradient-identity}
\end{align}
The first line gives \eqref{env:background-gradient}, since both
coefficients have positive lower bounds on $\overline P$.
\end{proof}

For a smooth function $f$ and a smooth metric $h$, write
\[
 a=\nabla_h f,\qquad H=\nabla_h^2f,\qquad
 B=\sqrt{2+|a|_h^2}.
\]
We call $f$ a \emph{strict envelope} for $h$ on a set if $a\ne0$
there and
\begin{equation}\label{env:condition}
 H(a,a)+(2+|a|_h^2)
       \max_{\substack{|e|_h=1\\df(e)=0}}H(e,e)>0.
\end{equation}
The maximum exists because the transverse unit sphere is compact.
The condition is open in the second derivatives of $f$ and the first
derivatives of $h$, as long as $df\ne0$.

The function $f_0$ is a strict envelope for $g_0$ on $\overline P$.
Indeed, with $e_\theta=(\sin p)^{-1}\partial_\theta$,
\[
 a_0=\cot p\,\partial_p,\qquad
 H_0(\partial_p,\partial_p)=-\csc^2p,\qquad
 H_0(e_\theta,e_\theta)=\cot^2p,
\]
and the expression in \eqref{env:condition}, tested with $e_\theta$,
is
\[
 -\cot^2p\,\csc^2p+(2+\cot^2p)\cot^2p=\cot^2p>0.
\]
Fix a smooth-topology open neighborhood $\mathcal U$ of $g_0$,
sufficiently small in $C^1$ that $f_0$ remains strict for every
$h\in\mathcal U$ on $\overline P$. Also arrange uniform comparison
of these metrics with $g_0$ on $M$ and, for one fixed $K\ge1$,
\begin{equation}\label{env:metric-comparison}
 K^{-1}I\le (h_{ij}(x))\le KI
       \quad (x\in\overline P,\ h\in\mathcal U).
\end{equation}
The finite-stage backgrounds used below satisfy the additional condition
that $h-g_0$ has compact support in $P$.

\subsection{Amplifying the gradient}

We reserve one direction for the old gradient and one transverse direction
for the maximizing Hessian test in \eqref{env:condition}. Corrugation in
the remaining two directions supplies a large positive tangential Hessian
that compensates for its negative radial Hessian. When this tangential
term does not dominate, the reserved transverse direction retains the
original strictness.

\begin{lemma}[Uniform amplification]\label{env:amplify}
There is $d_*>0$, depending only on the comparison constant $K$ and
fixed auxiliary functions, with the following property. Suppose $h$ is
smooth, satisfies \eqref{env:metric-comparison} near $\overline D$,
and $f$ is a strict envelope for $h$ on a neighborhood of $\overline D$.
For every $b>0$ there is a smooth strict envelope $\widetilde f$ such
that
\begin{gather}
 \operatorname{supp}(\widetilde f-f)\Subset D,\qquad
 \|\widetilde f-f\|_{L^\infty}<b,\label{env:amplification-smallness}\\
 \int_D|\nabla_h\widetilde f|_h\,dx
 \ge(1+d_*)\int_D|\nabla_h f|_h\,dx.
 \label{env:amplification-gain}
\end{gather}
Consequently the gradient integral can be made arbitrarily large by a
perturbation compactly supported in $D$, with any prescribed positive
total sup-norm budget. No derivative bound on the resulting envelope
is asserted.
\end{lemma}

\begin{proof}
Fix $0<r_0<r_1<R<1/2$ and a nonzero nonnegative function
$s\in C_c^\infty((r_0,r_1))$. In the disk of radius $R$ about each
point of $\mathbb Z^2$, define
\[
 q(r)=-\int_r^R s(t)\,dt,
\]
and set $q=0$ off the disks. The resulting function on $\mathbb R^2$
is smooth and $\mathbb Z^2$-periodic: it is a negative constant near
each disk center and zero near each disk boundary. Where its derivatives
are nonzero, let $n,t$ be the radial and tangential Euclidean unit
vectors. Then
\begin{equation}\label{env:radial-bowl}
 Dq=s\,n^\flat,\qquad
 D^2q=s'\,n^\flat\otimes n^\flat
             +\frac{s}{r}\,t^\flat\otimes t^\flat.
\end{equation}
Outside those annuli all derivatives in \eqref{env:radial-bowl} vanish.
Fix also $0\le\chi\in C_c^\infty((-1,1)^4)$, equal to one on a
smaller concentric box. Nonnegativity and bounded second derivatives
give
\begin{equation}\label{env:bump-estimates}
 |s'|\le C\sqrt{s},\qquad |D\chi|\le C\sqrt{\chi}.
\end{equation}
For example, for a nonnegative one-variable function $v$ with
$\|v''\|_\infty\le L$, Taylor's inequality at displacement
$-v'/L$ gives $|v'|^2\le2Lv$; applying this along lines proves the
second estimate as well. Let
\[
 K_q=\sup s(s')_-,\qquad (s')_-=\max\{-s',0\},
\]
and fix once and for all $0<\delta\le1$ so small that
\begin{equation}\label{env:fixed-delta}
 2K_q\delta^2\le\frac1{8R}.
\end{equation}
This choice is independent of $f$, $h$, and all subsequent amplifications.

\paragraph{The perturbation and its errors.}
Subdivide $D$ into finitely many congruent coordinate cubes $C$ of
side $L$. At the center $x_C$ of each cube choose constant linear
coordinates $y$ centered at $x_C$ such that $h(x_C)=I$,
\[
 a(x_C)=A_Ce_1,
\]
and $e_2$ realizes the transverse maximum in \eqref{env:condition}
at $x_C$. Complete these to an $h(x_C)$-orthonormal frame. By
\eqref{env:metric-comparison}, one fixed $c>0$ ensures that the
$y$-box $|y_j|<\eta$, with $\eta=cL$, lies strictly inside $C$;
its coordinate volume is a fixed positive fraction of $|C|$, uniformly
over all cells. These frames need not vary continuously between cells.

Put $\epsilon=\eta^3$ and, in this box, add
\begin{equation}\label{env:corrugation}
 v_C(y)=\delta A_C\epsilon\,
       \chi(y/\eta)q\bigl((y_3,y_4)/\epsilon\bigr).
\end{equation}
Extend $v_C$ by zero and set $\widetilde f=f+\sum_Cv_C$. The sum
is finite, its terms have disjoint supports, and its support is compact
in $D$. All constants in the error estimates that follow may depend
on the current fixed $f,h$, but not on the cell or its fineness.
In particular, $A_C$ has a positive lower bound and a finite upper
bound. For brevity write $A=A_C$ in a given cell. Uniformly there,
$h=I+O(\eta)$ and
\begin{align}
 d\widetilde f
 &=A\bigl(e_1^\flat+\delta\chi s\,n^\flat+O(\eta)\bigr),
 \notag\\
 a_*:=\nabla_h\widetilde f
 &=A\bigl(e_1+\delta\chi s\,n+O(\eta)\bigr),
 \label{env:gradient-expansion}\\
 H_*:=\nabla_h^2\widetilde f
 &=H+\delta A(Q+\mathcal E),\qquad
 Q=\frac{\chi}{\epsilon}D^2q.
 \label{env:hessian-expansion}
\end{align}
Here $\chi$ and $q$ are evaluated at their arguments in
\eqref{env:corrugation}; $n,t$ lie in the $(y_3,y_4)$-plane.
The differential of the cutoff term is $O(\epsilon/\eta)$.
Two differentiations of \eqref{env:corrugation}, including the
connection term in the Hessian, give
\begin{equation}\label{env:error-raw}
 |\mathcal E|\le C\left(
  \frac{\sqrt\chi\,s}{\eta}+\frac{\epsilon}{\eta^2}
              +\chi s+\frac{\epsilon}{\eta}\right).
\end{equation}
The second term includes the constant central part of the negative
bowl. In particular it cannot be omitted when $Dq=0$.

Define nonnegative quantities, equal to zero off the radial annuli,
\[
 T=\frac{\chi s}{\epsilon},\qquad
 N=\frac{\chi(s')_-}{\epsilon}.
\]
The following estimates hold uniformly, including at the cutoff and
annulus boundaries:
\begin{align}
 |\mathcal E|+\eta^2N
 &\le C\eta^{1/2}(1+T),\label{env:error-uniform}\\
 N(\chi s)^2
 &\le\|s'\|_\infty\epsilon T^2,
 \qquad N(\chi s)^2\le K_qT.
 \label{env:radial-error}
\end{align}
Indeed,
\[
 \frac{\sqrt\chi\,s}{\eta}
   =\eta^{1/2}\sqrt{sT}\le C\eta^{1/2}\sqrt T,
 \qquad
 \eta^2N\le C\eta^{1/2}\sqrt{\chi T},
\]
while the other terms in \eqref{env:error-raw} are
$\eta,\eta^3T,\eta^2$. Finally,
$N(\chi s)^2=T\chi^2s(s')_-$, which proves both estimates in
\eqref{env:radial-error}. Equation \eqref{env:gradient-expansion}
also gives $d\widetilde f(e_1)=A(1+O(\eta))$; thus the new
gradient never vanishes once the cells are sufficiently small.

\paragraph{Strictness when $T$ is large.}
At a point with $T>0$, take the \emph{exactly} transverse vector
\begin{equation}\label{env:angular-test}
 v=t-\frac{d\widetilde f(t)}{d\widetilde f(e_1)}e_1,
 \qquad e=\frac{v}{|v|_h}.
\end{equation}
Since $n^\flat(t)=0$, this is $t+O(\eta)e_1$ before normalization.
Its radial component is exactly zero, and normalization in the variable
metric is scalar multiplication. Hence
\[
 Q(e,e)=\frac{T}{r|v|_h^2}\ge\frac{T}{2R}
\]
for all sufficiently small cells. By
\eqref{env:gradient-expansion}, the radial component of $a_*/A$ is
$\delta\chi s+O(\eta)$. Discarding the nonnegative tangential
part of $Q$ and using \eqref{env:radial-error} gives
\begin{equation}\label{env:gradient-fast-lower}
 \begin{split}
 Q(a_*,a_*)&\ge
 -A^2\bigl(2\delta^2N(\chi s)^2+C\eta^2N\bigr)\\
 &\ge-A^2\bigl(2\delta^2K_qT+C\eta^2N\bigr).
 \end{split}
\end{equation}
Also $|a_*|_h^2\ge A^2/2$. For $T\ge1$,
\eqref{env:error-uniform} therefore bounds the contribution of
$\delta A(Q+\mathcal E)$ to the strictness expression from below by
\[
 \delta A^3\left(\frac1{4R}-2K_q\delta^2
                         -C\eta^{1/2}\right)T.
\]
Here and below the constant $C$ may change with the current envelope;
the positive lower bound on $A$ absorbs the factors $1+A^{-2}$ from
the error terms. By \eqref{env:fixed-delta}, sufficiently small
$\eta$ makes this at least $\delta A^3T/(16R)$. The contribution
of the old Hessian $H$ is bounded below by $-C_H$ uniformly in the
cell and its point. Choose a finite threshold
\[
 T_0>\max\left\{1,\frac{32RC_H}{\delta\min_{\overline D}|a|_h^3}\right\}.
\]
Then \eqref{env:condition} holds at every point where $T\ge T_0$,
once the subdivision is sufficiently fine. Only $T_0$ and this
fineness threshold depend on the current Hessian and gradient bounds.

\paragraph{Strictness when $T$ is bounded.}
For $0\le T\le T_0$, the radial coefficient $N$ may still be large.
We use the reserved transverse direction and bound the negative
contribution to $Q(a_*,a_*)$ separately. Set
\begin{equation}\label{env:old-test}
 v=e_2-\frac{d\widetilde f(e_2)}{d\widetilde f(e_1)}e_1,
 \qquad e=\frac{v}{|v|_h}.
\end{equation}
This vector is exactly transverse, is $e_2+O(\eta)$ after
normalization, and has no fast-plane component. Consequently
$Q(e,e)=0$ exactly. Moreover $\chi s=\epsilon T\to0$ uniformly
in this regime, so $a_*=Ae_1+o(1)$. Equations
\eqref{env:gradient-fast-lower} and \eqref{env:radial-error}, using
the stronger bounded-$T$ estimate before replacing it by $K_qT$,
give
\[
 Q(a_*,a_*)\ge
 -A^2\bigl(2\delta^2\|s'\|_\infty\epsilon T_0^2
                  +C\eta^2N\bigr)=o(1).
\]
Also $\mathcal E=o(1)$ uniformly by
\eqref{env:error-uniform}. The old-Hessian contribution, tested with
\eqref{env:old-test}, converges uniformly to its value at the cell
center, where $e_2$ was chosen to maximize it. Those center values
have a common positive lower bound, by strictness of $f$ on
$\overline D$. Thus the full expression is positive for sufficiently
small $\eta$. The two regimes cover every point of the perturbation;
outside its support the original envelope is unchanged.

\paragraph{A fixed gain and a small value change.}
In every cell, \eqref{env:gradient-expansion} implies
\begin{equation}\label{env:norm-gain}
 \frac{|a_*|_h}{A}
   =\sqrt{1+\delta^2\chi^2s^2}+O(\eta),
 \qquad \frac{|a|_h}{A}=1+O(\eta).
\end{equation}
Choose $s_0>0$ such that the set $s\ge s_0$ has positive area in
one fast-variable period. Since $\epsilon/\eta=\eta^2\to0$,
periodicity shows that a fixed fraction $\beta>0$ of each whole
cell has both $\chi=1$ and $s\ge s_0$. The constant $\beta$
depends only on the fixed cutoffs and metric comparison: the central
cutoff box contains many fast periods, and the linear coordinate
Jacobians and box-to-cell volume ratios have uniform bounds. Set
\[
 \gamma=\sqrt{1+\delta^2s_0^2}-1>0.
\]
Integrating \eqref{env:norm-gain} on a cell yields
\[
 \int_C|a_*|_h\,dx
 \ge\int_C|a|_h\,dx+A(\beta\gamma-C\eta)|C|.
\]
For sufficiently small $\eta$, the last term is at least
$A\beta\gamma|C|/2$, while $\int_C|a|_h\,dx\le2A|C|$.
Thus \eqref{env:amplification-gain} follows with
$d_*=\beta\gamma/4$, independently of the current strictness margin
and derivative bounds.

Finally, disjointness of the supports gives
\[
 \|\widetilde f-f\|_\infty
 \le\delta\bigl(\max_C A_C\bigr)\epsilon\|q\|_\infty,
\]
which tends to zero with the cell size. This proves the one-step
assertion with any budget $b$. For any desired finite integral, a
finite number of gains by $1+d_*$ suffices. Assign these finitely
many steps portions of the total sup-norm budget and choose each
subdivision after fixing its current envelope. Their supports form a
finite union compactly contained in $D$. Strictness and nonvanishing
hold after every step. This is a finite construction of a smooth
envelope, with no limit of envelopes required.
\end{proof}

\subsection{Separated envelope and background regions}

\begin{lemma}[Gaps and arbitrary growth]\label{env:gaps}
Fix $h\in\mathcal U$ with $\operatorname{supp}(h-g_0)\Subset P$.
There are nested open coordinate boxes
\[
 \overline D\cup\operatorname{supp}(h-g_0)\subset U_0
 \Subset U_1\Subset U_2\Subset U_3\Subset U_4\Subset P.
\]
For every $L_*>0$ there is a smooth strict envelope $f$ for $h$ on
a neighborhood of $\overline P$, and a number $\sigma>0$, such that
\begin{gather}
 f>f_0+\sigma\quad\hbox{on }\overline U_0,
 \qquad
 f<f_0-\sigma\quad\hbox{on }\overline P\setminus U_1,
 \label{env:gap-inequalities}\\
 \int_D\sqrt{2+|\nabla_h f|_h^2}\,dx>L_*.
 \label{env:large-integral}
\end{gather}
The envelope and all these patches are fixed before the frequency $k$
is chosen.
\end{lemma}

\begin{proof}
The indicated compact set lies inside the coordinate cube $P$, so
nested coordinate boxes with these strict inclusions exist. Choose
$\zeta\in C_c^\infty(U_1)$ equal to one on a neighborhood of
$\overline U_0$. For a sufficiently small fixed $\tau>0$,
\[
 f^{(1)}=f_0+\tau(2\zeta-1)
\]
is still strict on a neighborhood of $\overline P$, by the
$C^2$-openness of \eqref{env:condition}. Its gaps relative to $f_0$
are exactly $\tau$ on $\overline U_0$ and $-\tau$ outside $U_1$.
Apply Lemma~\ref{env:amplify} finitely many times in $D$, with total
sup-norm change less than $\tau/4$, to make
$\int_D|\nabla_hf|_h\,dx>L_*$. The resulting $f$ is unchanged
outside a compact subset of $D$ and remains strict everywhere in
question. Taking $\sigma=\tau/2$ gives
\eqref{env:gap-inequalities}; the inequality $B\ge|\nabla_hf|_h$
gives \eqref{env:large-integral}.
\end{proof}
